\begin{frame}{log-odds 직관: 시그모이드는 ``대수적 기울기''}
  시그모이드를 ``반전''하면 이 변환이 왜 자연스러운지 선명해진다.
  $p = \sigma(z)$ 를 $z$ 에 대해 정리:
  \[
  z = \log \frac{p}{1-p} = \log \text{odds}(p)
  \]
  \begin{itemize}
    \item $\text{odds}(p) = p/(1-p)$ = \kb{승산(odds)} --- ``양성일
      가능성 : 음성일 가능성''의 비율.
    \item 시그모이드의 역함수가 \kb{log-odds}(로지트, logit)이므로,
      로지스틱회귀는 실제로 \textbf{log-odds를 선형 모델로
      설명}:
      \[\log \frac{h_{\boldsymbol{w}}(\boldsymbol{x})}{1-h_{\boldsymbol{w}}(\boldsymbol{x})} = \boldsymbol{w}^T\boldsymbol{x}\]
  \end{itemize}
  \vskip0.4em
  숫자: $\boldsymbol{w}^T\boldsymbol{x} = 1.0986$ 이면 $h_{\boldsymbol{w}}(\boldsymbol{x}) = 0.75$ --- 승산 $0.75/0.25 = 3$ 이
 므로 ``양성일 가능성은 음성의 3배''.
\end{frame}
